% ======================================================================
% SECTION 13: CONCLUSIONS
% Summary of accomplishments and persistent themes. 4 to 6 paragraphs.
% ======================================================================

[PROCESSING INSTRUCTION (Conclusions, 4 to 6 paragraphs):

(13.1) One paragraph that lists the seven proposed bills by short
title and the prior law each adapts or revises:
- HR 4501 Patient Self-Selection (Adaption of 21 CFR Part 50, FDA
  DCT Final Guidance 2024, 42 USC 300gg-8)
- HR 4502 Robot-and-Humanoid Choice (Revision of CA AB 2847, FDA AI
  Regulatory Decision-Making Draft 2025)
- HR 4503 Procedural Modification Authority (Adaption of OHRP Broad
  Consent 2017, 21st Century Cures Act PL 114-255)
- HR 4504 Doctor and Nurse Error Rate Reduction (Adaption of HTI-1
  DSI Final Rule 2023, FDA AI Decision-Making Draft 2025)
- HR 4505 Real-Time Patient-Sponsor Direct Communication (Revision
  of FDA RTCT Press Announcement April 2026, FDA DCT Guidance 2024)
- HR 4506 American Leadership in Medical AI and Robotics (New
  Federal Authorizing Statute Adapting FDORA Section 3209, 21st
  Century Cures Act PL 114-255)
- HR 4507 Patient Health Data Self-Custody (Revision of HHS HIPAA
  Right-to-Access 2025, ONC Cures Act Final Rule).

(13.2) One paragraph on the central thesis: cancer patients in U.S.
oncology clinical trials are the priority participants and must hold
statutory control over discovery, scheduling, robot selection,
procedural modification, real-time data feedback, error-rate
information, and self-custody of their health data. Advanced AI and
advanced robotics now offer the patient the opportunity to take that
control directly; the seven bills supply the legal infrastructure
that makes the control enforceable and measurable. Cite
\cite{kawchak_2026_19994945} and \cite{kawchak_2026_19119939}.

(13.3) One paragraph on the persistent themes of the paper:
- Adaption / Revision framing: every bill explicitly states the
  prior U.S. legislative document it builds from.
- Respectful FDA framing: every bill positions the FDA and other
  governing bodies as enabling rather than gating, while remaining
  opportunistic under the patient's control based primarily on
  their own interests.
- Individual patient + broad adoption framing: every bill names
  per-patient examples (PAT-2026-0042 from patient-journey/) AND
  broad-adoption deadlines (12, 18, 24 months across CMS-billable
  sites and 50 states).
- United States Number 1: every bill closes with the leadership
  framing required to keep the United States Number 1 in the world
  for patient benefit during the medical AI and robotics revolution.

(13.4) One paragraph on the implications for cancer patient control:
LLM-scale repository context (1M tokens) demonstrated in
new-trial/national-24-7-trial/paper/full-paper/, the four LLM
simulations summarized at \cite{kawchak_2026_19994945}, the 168-hour
sponsor extension at sponsor/final\_paper/168\_hours/, the 10-stage
single-patient journey at patient-journey/, and the 10-robot-category
patient instructions at patients/patient\_robot\_instructions\_fixed.tex
collectively show that the technical substrate is ready. The seven
bills supply the missing legal substrate. Together they enable the
cancer patient to take control of their disease in ways that prior
error-prone and physically and intellectually limited human-only care
could not support.

(13.5) One paragraph that restates the key limitations: no
congressional review, no Real-patient-cohort outcome study, draft
language only, vendor neutrality preserved. State that the work is a
legislative-design and operational demonstration, not a clinical
demonstration.

(13.6) One paragraph that points forward: the next Claude Code Opus
4.7 1M Max generation pass will populate this template into a 70+
page paper with full prose, color diagrams, and additional individual
patient and robot examples drawn from the AVAILABLE DIRECTORIES; the
author will then perform the senior-author white-space and table-
column-width formatting pass to remove any orphans, widows, or large
empty page regions before publication. Cite the future Zenodo
deposition at \cite{kawchak_2026_20045457}, the FDA RTCT announcement
at \cite{fda_rtct_press_2026}, the four-simulation paper at
\cite{kawchak_2026_19994945}, the sponsor paper at
\cite{kawchak_2026_19396256}, the National Platform paper at
\cite{kawchak_2026_19244918}, and the patient instructions paper at
\cite{kawchak_2026_18810541}.]
